\documentclass[11pt,a4paper]{article}
\usepackage[a4paper,margin=2.2cm]{geometry}
\usepackage[utf8]{inputenc}
\usepackage[T1]{fontenc}
\usepackage[ngerman]{babel}
\usepackage{lmodern}
\usepackage{amsmath,amssymb,mathtools}
\usepackage{microtype}
\usepackage{xcolor}
\usepackage{tcolorbox}
\tcbuselibrary{breakable,skins}
\usepackage{tikz}
\usetikzlibrary{arrows.meta,calc,decorations.markings,patterns,positioning}
\usepackage{pgfplots}
\pgfplotsset{compat=1.18}
\usepackage{enumitem}
\usepackage{array,booktabs,longtable}
\usepackage{hyperref}
\hypersetup{colorlinks=true,linkcolor=blue!60!black,urlcolor=blue!60!black}
\setcounter{secnumdepth}{0}

\newcommand{\C}{\mathbb C}
\newcommand{\R}{\mathbb R}
\newcommand{\N}{\mathbb N}
\newcommand{\Z}{\mathbb Z}
\newcommand{\ii}{\mathrm i}
\newcommand{\ee}{\mathrm e}
\newcommand{\dd}{\,\mathrm d}
\newcommand{\Arg}{\operatorname{Arg}}
\newcommand{\Log}{\operatorname{Log}}
\newcommand{\Res}{\operatorname{Res}}
\newcommand{\Ree}{\operatorname{Re}}
\newcommand{\Imm}{\operatorname{Im}}

\newtcolorbox{importantbox}[1][]{breakable, enhanced, colframe=red!75!black, colback=red!2, coltitle=black, fonttitle=\bfseries, title=Wichtig, #1}
\newtcolorbox{fragebox}[1][]{breakable, enhanced, colframe=black!70, colback=gray!3, fonttitle=\bfseries, title=Frage, #1}
\newtcolorbox{antwortbox}[1][]{breakable, enhanced, colframe=blue!60!black, colback=blue!2, fonttitle=\bfseries, title=Antwort, #1}
\newtcolorbox{tippbox}[1][]{breakable, enhanced, colframe=purple!70!black, colback=purple!3, fonttitle=\bfseries, title=Tipp, #1}
\newtcolorbox{kochrezeptbox}[1][]{breakable, enhanced, colframe=green!45!black, colback=green!2, fonttitle=\bfseries, title=Kochrezept, #1}
\newtcolorbox{beispielbox}[1][]{breakable, enhanced, colframe=black!65, colback=white, fonttitle=\bfseries, title=Beispiel, #1}

\tikzset{axis/.style={->,thick}, vector/.style={->,very thick,blue!70!black}, patharrow/.style={postaction={decorate},decoration={markings,mark=at position .55 with {\arrow{Stealth}}}}}
\newcommand{\complexplane}[4]{%
\begin{tikzpicture}[scale=#1]
  \draw[axis] (-#2,0)--(#2,0) node[right] {$\Re(z)$};
  \draw[axis] (0,-#3)--(0,#3) node[above] {$\Im(z)$};
  \draw[step=1,gray!20,thin] (-#2,-#3) grid (#2,#3);
  #4
\end{tikzpicture}}

\title{Mathematische Methoden\\Übungsskript nach Wochen}
\author{}
\date{Woche 1 bis Woche 14}

\begin{document}
\maketitle
\tableofcontents
\newpage

\section*{Vorbemerkung}
Dieses Skript fasst den Übungsstoff der Vorlesung Mathematische Methoden zusammen. Ich kann keine Garantie für die Korrektheit und Vollständigkeit des Skripts geben. Falls euch
Fehler im Skript auffallen, schreibt mir gerne eine E-Mail an rjahnke@ethz.ch. Ich werde die
aktuellste Version des Skriptes immer auf meine Website hochladen.

% =========================================================
\section{Woche 1: Komplexe Zahlen}

\subsection{1.1 Definition der komplexen Zahlen}
Die komplexen Zahlen können als reeller Vektorraum mit einer speziellen Multiplikation verstanden werden:
\[
\C := \left(\R^2,+,\cdot\right), \qquad
\binom{x}{y}\cdot\binom{x'}{y'}=
\binom{xx'-yy'}{xy'+x'y}.
\]
Komplexe Zahlen sind also zunächst Vektoren in $\R^2$. Neu ist nicht die Addition, sondern die Multiplikation.

\begin{importantbox}
Komplexe Zahlen sind ein reeller Vektorraum mit einer speziellen Multiplikation. Die Addition funktioniert wie in $\R^2$; die komplexe Multiplikation ist der wesentliche neue Inhalt.
\end{importantbox}

\subsection{1.2 Normalform und Multiplikation}
Da eine komplexe Zahl $z$ als reeller Vektor aufgefasst werden kann, besteht sie aus zwei reellen Zahlen $x,y\in\R$:
\[
 z=x+\ii y, \qquad \ii^2=-1.
\]
Hier kennzeichnet $\ii$ die zweite Komponente. Diese Darstellung heißt \emph{Normalform}. Dabei ist $x$ die erste Komponente und $y$ die zweite Komponente.

Die spezielle Multiplikation stimmt mit der Regel $\ii^2=-1$ überein:
\[
(x+\ii y)(x'+\ii y')=xx'+\ii xy'+\ii x'y+\ii^2yy'
=(xx'-yy')+\ii(xy'+x'y).
\]
Die Addition passt ebenfalls zur Vektoraddition:
\[
(x+\ii y)+(x'+\ii y')=(x+x')+\ii(y+y').
\]

\subsection{1.3 Darstellung von $\ii$ als Vektor}
\begin{fragebox}
Wie würde man $\ii$ als Vektor darstellen?
\end{fragebox}
\begin{antwortbox}
\[
\ii=0+\ii\cdot 1 \quad\Longleftrightarrow\quad \ii=\binom{0}{1}.
\]
\end{antwortbox}

\subsection{1.4 Realteil, Imaginärteil und komplexe Ebene}
Die erste Komponente bezeichnet man als Realteil:
\[
\Ree(z)=\Ree(x+\ii y)=x.
\]
Analog heißt die zweite Komponente Imaginärteil:
\[
\Imm(z)=y.
\]
Komplexe Zahlen lassen sich daher in der Ebene darstellen.

\begin{center}
\complexplane{0.75}{3.3}{3.3}{
  \draw[vector] (0,0)--(2,2) node[above right] {$z=2+2\ii$};
  \draw[dashed] (2,0)--(2,2)--(0,2);
  \node[below] at (2,0) {$x$};
  \node[left] at (0,2) {$y$};
}
\end{center}

\begin{importantbox}
Die komplexe Ebene ist als Vektorraum wie $\R^2$, aber statt des gewöhnlichen Skalarprodukts benutzen wir die komplexe Multiplikation.
\end{importantbox}

Die komplexe Multiplikation ermöglicht Rotation und Streckung in der Ebene. Multipliziert man zwei komplexe Zahlen als Vektoren, dann wird geometrisch ein Vektor rotiert und gestreckt.

\subsection{1.5 Multiplikation als Rotation und Streckung}
\begin{beispielbox}
Wir stellen $z=1+2\ii$ als Vektor dar und multiplizieren mit $\ii$:
\[
z\cdot\ii=(1+2\ii)\ii=\ii+2\ii^2=-2+\ii.
\]
\end{beispielbox}

\begin{center}
\begin{tikzpicture}[scale=0.75]
  \draw[axis] (-3.2,0)--(3.2,0) node[right] {$\Re(z)$};
  \draw[axis] (0,-3.2)--(0,3.2) node[above] {$\Im(z)$};
  \draw[step=1,gray!20,thin] (-3,-3) grid (3,3);
  \draw[vector] (0,0)--(1,2) node[above right] {$1+2\ii$};
  \draw[vector,red!70!black] (0,0)--(-2,1) node[above left] {$-2+\ii$};
  \draw[->,thick] (0.8,1.6) arc (63:153:1.8);
  \node at (-0.3,2.35) {$90^\circ$};
\end{tikzpicture}
\end{center}

Zur Streckung verwendet man den Betrag:
\[
 |z|=|x+\ii y|=\sqrt{x^2+y^2}.
\]
Beispiel: $|1+2\ii|=\sqrt{5}$ und $|-2+\ii|=\sqrt{5}$. Die Streckung ist also $1$.

Zur Rotation verwendet man das Argument. Für $z=x+\ii y\neq 0$ ist der Hauptwert definiert durch
\begin{importantbox}[title=Definition des Arguments]
\begin{minipage}[c]{0.68\linewidth}
\[
\Arg(z):=
\begin{cases}
\displaystyle \arccos\left(\frac{x}{|z|}\right), & y\geq 0,\\[1.2ex]
\displaystyle -\arccos\left(\frac{x}{|z|}\right), & y<0.
\end{cases}
\]
\[
\Arg(z)\in(-\pi,\pi],
\qquad
\Arg(x+\ii y)=\operatorname{atan2}(y,x).
\]
\end{minipage}\hfill
\begin{minipage}[c]{0.27\linewidth}
\centering
\begin{tikzpicture}[scale=0.9]
  \draw[axis] (-1.55,0)--(1.55,0) node[right] {$\Re(z)$};
  \draw[axis] (0,-1.35)--(0,1.45) node[above] {$\Im(z)$};
  \draw[blue!65!black,very thick,->]
    (1.05,0) arc[start angle=0,end angle=170,radius=1.05];
  \draw[blue!65!black,very thick,->]
    (1.05,0) arc[start angle=0,end angle=-170,radius=1.05];
  \node[above left] at (-0.82,0.35) {$+\pi$};
  \node[below left] at (-0.82,-0.35) {$-\pi$};
\end{tikzpicture}
\end{minipage}
\end{importantbox}
Das Argument ist also der eindeutig gewählte Winkel des Vektors $(x,y)$ zur positiven reellen Achse.

Im Beispiel gilt
\[
\Arg(1+2\ii)=\arccos\left(\frac1{\sqrt5}\right)\approx 63{,}435^\circ,
\]
\[
\Arg(-2+\ii)=\arccos\left(-\frac2{\sqrt5}\right)\approx 153{,}435^\circ.
\]
Die Rotation ist also $153{,}435^\circ-63{,}435^\circ=90^\circ$.

\begin{fragebox}
War das ein Spezialfall oder ergibt eine Multiplikation mit $\ii$ immer eine reine $90^\circ$-Rotation? Warum?
\end{fragebox}
\begin{antwortbox}
Für $z\neq0$ ist
\[
 z\ii = (x+\ii y)\ii=-y+\ii x.
\]
Der Vektor $(x,y)$ wird zu $(-y,x)$ gedreht. Das ist genau eine Drehung um $90^\circ$ gegen den Uhrzeigersinn.
\end{antwortbox}

\subsection{1.6 Polarform}
Neben der Normalform gibt es die Polarform:
\[
z=x+\ii y=r(\cos\varphi+\ii\sin\varphi),\qquad r\in\R_+,\ \varphi\in[-\pi,\pi].
\]
Aus der Euler-Formel folgt
\[
\ee^{\ii\varphi}=\cos\varphi+\ii\sin\varphi,
\]
also
\begin{importantbox}
\[
 z=r\ee^{\ii\varphi}.
\]
Die Polarform eignet sich besonders gut für Multiplikation, Division und Potenzen.
\end{importantbox}

\subsection{1.7 Bestimmung von $r$ und $\varphi$}
\begin{fragebox}
Was ist $r$ und $\varphi$ für $z=\ii$?
\end{fragebox}
\begin{antwortbox}
\[
r=|\ii|=1,\qquad \varphi=\frac\pi2,
\]
also
\[
\ii=\ee^{\ii\pi/2}.
\]
\end{antwortbox}

Seien $z_1=r_1\ee^{\ii\varphi}$ und $z_2=r_2\ee^{\ii\psi}$. Dann gilt
\[
 z_1z_2=r_1r_2\ee^{\ii(\varphi+\psi)},
\qquad
 \frac{z_1}{z_2}=\frac{r_1}{r_2}\ee^{\ii(\varphi-
\psi)}.
\]
Damit ist bewiesen: Die Beträge multiplizieren sich, die Winkel addieren sich.

\subsection{1.8 Rechnen mit Polarform}
Die Normalform ist nützlich für Addition und Subtraktion. Die Polarform ist nützlich für Multiplikation, Division und Potenzen.

\begin{fragebox}
Was ergibt $\ii^{202}$?
\end{fragebox}
\begin{tippbox}
Verwende die Polarform $\ii=\ee^{\ii\pi/2}$.
\end{tippbox}
\begin{antwortbox}
\[
\ii^{202}=\left(\ee^{\ii\pi/2}\right)^{202}=\ee^{\ii 101\pi}=-1.
\]
\end{antwortbox}

\subsection{1.9 Konjugation}
Wir definieren die Konjugierte einer komplexen Zahl $z=x+\ii y$ als:
\[
 \overline z=x-\ii y,\qquad z\overline z=|z|^2=x^2+y^2.
\]
Damit kann man beispielsweise schreiben:
\[
 \frac{1}{z}=\frac{\overline z}{z\overline z}=\frac{x-\ii y}{x^2+y^2}.
\]

\begin{fragebox}
Was ist die Konjugierte von $\ee^{\ii\varphi}$?
\end{fragebox}
\begin{antwortbox}
\[
\overline{\ee^{\ii\varphi}}=\ee^{-\ii\varphi}.
\]
\end{antwortbox}

\subsection{2.0 Zusammenhang zwischen Realteil und Konjugation}
\begin{importantbox}
\[
\Ree(z)=\frac12(z+\overline z).
\]
\end{importantbox}

\begin{fragebox}
Was ergibt sich analog für $\Imm(z)$ in Abhängigkeit von $z$ und $\overline z$?
\end{fragebox}
\begin{antwortbox}
\[
\Imm(z)=\frac{z-\overline z}{2\ii}.
\]
\end{antwortbox}

\subsection{2.1 Normform und Konjugation}
\begin{fragebox}
Gilt
\[
 \Imm\bigl((-5+3\ii)(-7-5\ii)\bigr)=\Imm((-5+3\ii)(-7-5\ii))?
\]
\end{fragebox}
\begin{tippbox}
\[
\Imm(\overline z)=-\Imm(z).
\]
\end{tippbox}
\begin{antwortbox}
Ja
\end{antwortbox}

\begin{fragebox}
Berechne die Normform von
\[
z=\frac{27}{4+3\ii}.
\]
\end{fragebox}
\begin{antwortbox}
\[
 z=\frac{27}{4+3\ii}\cdot\frac{4-3\ii}{4-3\ii}
 =\frac{27(4-3\ii)}{4^2+3^2}
 =\frac{108}{25}-\frac{81}{25}\ii.
\]
\end{antwortbox}

\subsection{2.2 Rotation durch $\ee^{\ii\varphi}$}
Wir haben gesehen: Multiplikation mit $\ee^{\ii\varphi}$ bedeutet Rotation um den Winkel $\varphi$.

\begin{fragebox}
Was bedeutet $z\ii^p$?
\end{fragebox}
\begin{tippbox}
Schreibe $\ii=\ee^{\ii\pi/2}$ und benutze, dass Multiplikation mit $\ee^{\ii\varphi}$ eine Drehung um $\varphi$ ist.
\end{tippbox}
\begin{antwortbox}
Da $\ii^p=\ee^{\ii p\pi/2}$, bedeutet Multiplikation mit $\ii^p$ eine Rotation um $p\pi/2$.
\end{antwortbox}

\subsection{2.3 Nicht-Eindeutigkeit von $\ee^{\ii x}$}
\[
\ee^{2\pi\ii}=1=\ee^0.
\]
\begin{importantbox}
$\ee^{\ii x}$ ist nicht eindeutig im Argument, denn
\[
\ee^{\ii(\varphi+2\pi k)}=\ee^{\ii\varphi},\qquad k\in\Z.
\]
Daher ist $z$ in Polarform allgemein nicht eindeutig darstellbar, wenn man nicht einschränkt.
\end{importantbox}
Eine Einschränkung auf $(-\pi,\pi]$ ist nur unproblematisch, solange man keine Wurzeln, Potenzen und Logarithmen betrachtet.

\subsection{2.5 $n$-te Wurzeln}
Im Reellen ist Wurzelziehen nicht eindeutig: Es gibt etwa $\pm\sqrt a$. Analog gibt es im Komplexen $n$ verschiedene $n$-te Wurzeln.

Sei
\[
z=r\ee^{\ii(\varphi+2\pi k)}.
\]
Dann sind die $n$-ten Wurzeln
\begin{importantbox}
\[
w_k=\sqrt[n]{r}\,\ee^{\ii(\varphi+2\pi k)/n},\qquad k=0,1,\ldots,n-1.
\]
\end{importantbox}
Die Werte wiederholen sich periodisch in $k$.

\begin{fragebox}
Wo liegen alle $n$ Wurzeln?
\end{fragebox}
\begin{antwortbox}
Sie liegen auf einem Kreis um den Ursprung mit Radius $\sqrt[n]{r}$.
\end{antwortbox}

\begin{center}
\begin{tikzpicture}[scale=1]
  \draw[axis] (-1.5,0)--(1.7,0) node[right] {$\Re$};
  \draw[axis] (0,-1.5)--(0,1.7) node[above] {$\Im$};
  \draw (0,0) circle (1);
  \foreach \ang/\lab in {25/$k=0$,115/$k=1$,205/$k=2$,295/$k=3$}{
    \fill ({cos(\ang)},{sin(\ang)}) circle (1.5pt) node[shift={(\ang:0.28)}] {\lab};
  }
\end{tikzpicture}
\end{center}

% =========================================================
\section{Woche 2: Mehrdeutigkeit komplexer Zahlen}

\subsection{1.2 Mehrdeutigkeit der Polarform}
\[
\ee^{2\pi\ii}=1=\ee^{0}.
\]
\begin{importantbox}
Die Exponentialdarstellung $\ee^{\ii x}$ ist nicht eindeutig. Denn
\[
 z=r\ee^{\ii(\varphi+2\pi k)},\qquad k\in\Z,
\]
ist dieselbe komplexe Zahl für alle $k$.
\end{importantbox}
Daher ist die Polarform ohne Einschränkung nicht eindeutig.

\subsection{1.3 $n$-te Wurzeln}
Sei $w^n=z$, $z=r\ee^{\ii(\varphi+2\pi k)}$. Dann
\[
w_k=\sqrt[n]{r}\,\ee^{\ii(\varphi+2\pi k)/n},\qquad k=0,\ldots,n-1.
\]
Nur diese $n$ Werte sind verschieden; danach wiederholen sie sich.

\begin{fragebox}
Wo liegen alle $n$ Wurzeln?
\end{fragebox}
\begin{antwortbox}
Auf einem Kreis um den Ursprung mit Abstand $\sqrt[n]{r}$.
\end{antwortbox}

\subsection{1.4 Beispiel: fünfte Wurzeln von $\ii$}
\begin{beispielbox}
Gesucht sind die Lösungen von $z^5=\ii$. Weil
\[
\ii=\ee^{\ii\pi/2},
\]
gilt
\[
z_k=\ee^{\ii\left(\frac{\pi/2+2\pi k}{5}\right)},\qquad k=0,1,2,3,4.
\]
\end{beispielbox}

\subsection{1.5 Beispiel: Wurzeln von $z^5=3\ee^{3\pi\ii/2}$}
\begin{fragebox}
Was ist die Lösung von $z^5=3\ee^{3\pi\ii/2}$?
\end{fragebox}
\begin{tippbox}
Wandle die rechte Seite zuerst in Polarform um und verwende dann die Formel für fünfte Wurzeln.
\end{tippbox}
\begin{antwortbox}
\[
z_k=\sqrt[5]{3}\,\ee^{\ii\frac{3\pi/2+2\pi k}{5}},\qquad k=0,1,2,3,4.
\]
\end{antwortbox}

\subsection{1.6 Komplexer Logarithmus}
Ähnlich wie bei Wurzeln ist auch der komplexe Logarithmus nicht eindeutig.

\begin{importantbox}
Falls $\ee^w=z$, dann ist $w$ ein Logarithmus von $z$.
\end{importantbox}

Beispiel: $\ee^w=1$. Normalerweise wäre $w=\ln(1)=0$, aber im Komplexen gilt
\[
\ee^{2\pi\ii k}=1,
\]
also
\[
w=2\pi\ii k,\qquad k\in\Z.
\]
Es gibt unendlich viele Werte.

Allgemein schreiben wir $z$ in Polarform:
\[
z=r\ee^{\ii(\Arg(z)+2\pi k)}.
\]
Dann
\[
\log(z)=\ln|z|+\ii(\Arg(z)+2\pi k),\qquad k\in\Z.
\]
Für den Hauptwert setzt man $k=0$ und beschränkt das Argument.

\begin{importantbox}
\[
\Log(z)=\ln|z|+\ii\Arg(z),\qquad \Arg(z)\in(-\pi,\pi].
\]
\end{importantbox}

\subsection{1.7 Beispiele zum Logarithmus}
\begin{fragebox}
Was ist $\Log(-\sqrt2+\ii\sqrt2)$?
\end{fragebox}
\begin{tippbox}
Verwende die Polarform: zuerst Betrag und Hauptargument bestimmen.
\end{tippbox}
\begin{antwortbox}
Der Betrag ist $2$ und das Argument ist $3\pi/4$. Daher
\[
\Log(-\sqrt2+\ii\sqrt2)=\ln(2)+\ii\frac{3\pi}{4}.
\]
\end{antwortbox}

\begin{fragebox}
Was ist $\Log((-\sqrt2+\ii\sqrt2)^2)$?
\end{fragebox}
\begin{antwortbox}
Zunächst ist $(-\sqrt2+\ii\sqrt2)^2=-4\ii$. Damit ist
\[
\Log((-\sqrt2+\ii\sqrt2)^2)=\Log(-4\ii)=\ln(4)-\ii\frac{\pi}{2}.
\]
Achtung: $2\Log(z)=2\ln(2)+\ii\frac{3\pi}{2}=\ln(4)+\ii\frac{3\pi}{2}$ liegt nicht im Hauptzweig und muss durch Subtraktion von $2\pi\ii$ auf $\ln(4)-\ii\pi/2$ gebracht werden.
\end{antwortbox}

\begin{importantbox}
Die Regel $\Log(z^n)=n\Log(z)$ gilt im Hauptzweig nicht ohne Anpassung des Arguments auf $(-\pi,\pi]$.
\end{importantbox}

% =========================================================
\section{Woche 3: Fundamentalsatz}

\subsection{1.1 Beispiel: fünfte Wurzeln von $1$}
Für die fünften Wurzeln von $1$ gilt
\[
1=\ee^{2\pi\ii k},\qquad k\in\Z,
\]
und damit
\[
z_k=\ee^{2\pi\ii k/5},\qquad k=0,1,2,3,4.
\]

\begin{center}
\begin{tikzpicture}[scale=1]
  \draw[axis] (-1.4,0)--(1.6,0) node[right] {$\Re$};
  \draw[axis] (0,-1.4)--(0,1.6) node[above] {$\Im$};
  \draw[red!70!black] (0,0) circle (1);
  \foreach \k in {0,...,4}{
    \pgfmathsetmacro{\ang}{72*\k}
    \fill[red!70!black] ({cos(\ang)},{sin(\ang)}) circle (1.8pt) node[shift={(\ang:0.25)}] {$k=\k$};
  }
\end{tikzpicture}
\end{center}
Man kann die Wurzeln gegen den Uhrzeigersinn wählen, im Uhrzeigersinn oder in anderer Reihenfolge, etwa $k=0,1,-1,2,-2$.

\subsection{1.2 Konjugierte Wurzeln}
Wenn $z$ eine Wurzel von $z^5=1$ ist, dann ist auch $\overline z$ eine Lösung. Allgemeiner gilt:

\begin{importantbox}
Falls $z$ eine Wurzel einer reellen Zahl ist, also $z^n=a\in\R$, dann ist auch $\overline z$ eine Lösung.
\end{importantbox}

\subsection{1.3 Fundamentalsatz der Algebra}
Das Polynom $z^n-a$ hat $n$ Nullstellen, denn es hat $n$ Wurzeln. Das verallgemeinert der Fundamentalsatz der Algebra.

\begin{importantbox}
Sei
\[
p(z)=a_0+a_1z+\cdots+a_nz^n,
\qquad n\in\N,
\qquad a_n\neq0.
\]
Dann besitzt $p$ in $\C$ genau $n$ Nullstellen, wenn Vielfachheiten mitgezählt werden.
\end{importantbox}
Achtung: Die Nullstellen müssen nicht verschieden sein. Zum Beispiel ist $1$ eine doppelte Nullstelle von $(z-1)^2=0$.

\subsection{1.4 Frage zum Fundamentalsatz}
\begin{fragebox}
Sei
\[
p(z)=z^5-15z^4+95z^3-315z^2+544z-390.
\]
$3+2\ii$ ist eine Nullstelle. Ist auch $3-2\ii$ eine Nullstelle? Wie viele Nullstellen hat dieses Polynom?
\end{fragebox}
\begin{antwortbox}
Da das Polynom reelle Koeffizienten hat, tritt zu jeder nicht-reellen Nullstelle auch die konjugierte Nullstelle auf. Also ist $3-2\ii$ ebenfalls eine Nullstelle. Da der Grad $5$ ist, hat das Polynom insgesamt $5$ Nullstellen, mit Vielfachheiten gezählt.
\end{antwortbox}

\subsection{1.5 Verständnis der komplexen Ebene}
\begin{fragebox}
Zeichne
\[
\left\{z\in\C\,\middle|\,\left|\frac{z-1}{z+\ii}\right|=3\right\}.
\]
\end{fragebox}
\begin{antwortbox}
Setze $z=x+\ii y$. Dann
\[
|x-1+\ii y|^2=9|x+\ii(y+1)|^2.
\]
Also
\[
(x-1)^2+y^2=9(x^2+(y+1)^2).
\]
Ausmultiplizieren liefert
\[
x^2-2x+1+y^2=9x^2+9y^2+18y+9.
\]
Damit
\[
8x^2+2x+8y^2+18y+8=0.
\]
Durch quadratische Ergänzung erhält man
\[
\left(x+\frac18\right)^2+\left(y+\frac98\right)^2=\frac{81}{64}+\frac{1}{64}-1=\frac{18}{64}=\frac{9}{32}.
\]
Das ist ein Kreis mit Mittelpunkt $(-1/8,-9/8)$ und Radius $\sqrt{9/32}=\frac{3}{4\sqrt2}$.
\end{antwortbox}

\begin{center}
\begin{tikzpicture}[scale=1.2]
  \draw[axis] (-2.1,0)--(1.4,0) node[right] {$\Re$};
  \draw[axis] (0,-2.3)--(0,1.0) node[above] {$\Im$};
  \draw[blue!70!black,thick] (-0.125,-1.125) circle ({3/(4*sqrt(2))});
  \fill (-0.125,-1.125) circle (1.5pt) node[below left] {$(-\frac18,-\frac98)$};
\end{tikzpicture}
\end{center}

% =========================================================
\section{Woche 4: Komplexe Funktionen}

\subsection{1.1 Definition einer komplexen Funktion}
Eine komplexe Funktion ist eine Abbildung
\[
f:D\to\C,
\]
wobei $D$ entweder eine Teilmenge der reellen Zahlen oder eine Teilmenge der komplexen Zahlen sein kann.

\begin{beispielbox}
\[
f(x)=x+\ii \quad\text{mit reellem Definitionsbereich},
\]
\[
f(z)=z^2 \quad\text{mit komplexem Definitionsbereich}.
\]
\end{beispielbox}

\begin{importantbox}
Der Definitionsbereich $D$ kann eine Teilmenge von $\R$ oder eine Teilmenge von $\C$ sein. Das ist für Konvergenzfragen wichtig.
\end{importantbox}

\subsection{1.2 Warum ist $D$ wichtig?}
Wenn wir eine Funktion betrachten, möchten wir wissen, wo sie konvergiert. Wir wählen daher $D$ so, dass die Funktion auf $D$ konvergiert.

Viele Funktionen lassen sich als Potenzreihe schreiben:
\[
f(z)=\sum_{n=0}^{\infty}a_nz^n.
\]
Dann berechnet man den Konvergenzradius $R$ und wählt
\[
D=B(0,R),
\]
den offenen Ball um $(0,0)$ mit Radius $R$.

\subsection{1.3 Definitionsbereich einer verschobenen Potenzreihe}
\begin{fragebox}
Welches $D$ würden wir für
\[
f(z)=\sum_{n=0}^{\infty}a_n(z-z_0)^n
\]
wählen?
\end{fragebox}
\begin{antwortbox}
\[
D=B(z_0,R).
\]
\end{antwortbox}

\subsection{1.5 Nicht existierender Grenzwert}
\begin{fragebox}
Existiert
\[
\lim_{z\to0}\frac{\overline z}{z}?
\]
\end{fragebox}
\begin{antwortbox}
Nähert man sich auf der reellen Achse, also $z=x$, so gilt $\overline z/z=x/x=1$. Nähert man sich auf der imaginären Achse, also $z=\ii y$, so gilt
\[
\frac{\overline z}{z}=\frac{-\ii y}{\ii y}=-1.
\]
Die Grenzwerte stimmen nicht überein, also existiert der Grenzwert nicht.
\end{antwortbox}

\subsection{1.6 Komplexe Ableitung}
Ableitungen komplexer Funktionen funktionieren im Prinzip wie bei reellen Funktionen. Falls
\[
\lim_{z\to z_0}\frac{f(z)-f(z_0)}{z-z_0}
\]
existiert, ist $f$ in $z_0$ komplex differenzierbar.

Es gelten die üblichen Regeln: differenzierbar impliziert stetig, Quotientenregel, Produktregel, Kettenregel.

\subsection{1.7 Holomorphie}
\begin{importantbox}
\begin{enumerate}[label=\arabic*.]
\item $f$ ist holomorph auf $U$, falls $f$ auf ganz $U$ komplex differenzierbar ist.
\item $f$ ist holomorph in $z_0$, falls $f$ in einer offenen Menge um $z_0$ holomorph ist.
\end{enumerate}
\end{importantbox}

Achtung: Komplex differenzierbare Funktionen verhalten sich zwar formal wie reelle differenzierbare Funktionen, Holomorphie ist aber eine sehr starke Eigenschaft.

\subsection{1.8 Beispiele für Holomorphie}
\begin{fragebox}
Welche dieser Funktionen sind holomorph auf $D=B(0,1)$?
\begin{enumerate}[label=\roman*)]
\item $f(z)=\sum_{n=0}^{\infty}2^nz^n$
\item $f(z)=\frac{1}{1+3z}$
\item $f(z)=\ee^{z^2}$
\item $f(z)=\frac1{1-z\overline z}$
\item $f(z)=\frac1{1-z^2}$
\item $f(z)=1+z+z^2+z^3+\cdots$
\end{enumerate}
\end{fragebox}
\begin{tippbox}
Nicht stetig bedeutet nicht differenzierbar. Außerdem ist eine Abhängigkeit von $\overline z$ ein typischer Hinweis darauf, dass eine Funktion nicht holomorph ist.
\end{tippbox}
\begin{antwortbox}
Holomorph sind ii), iii), v) und vi), sofern die jeweilige Singularität nicht im Gebiet liegt bzw. die Reihe dort konvergiert. i) hat Konvergenzradius $1/2$ und ist daher nicht auf ganz $B(0,1)$ holomorph. iv) hängt von $\overline z$ ab und ist daher im Allgemeinen nicht holomorph. Stetigkeit allein ist nicht Differenzierbarkeit.
\end{antwortbox}

% =========================================================
\section{Woche 5: Holomorphie}

\subsection{1.1 Zerlegung in Real- und Imaginärteil}
Eine komplexe Funktion lässt sich schreiben als
\[
f(z)=f(x+\ii y)=u(x,y)+\ii v(x,y),
\]
wobei $u,v$ reelle Funktionen sind.

Wir betrachten drei Ableitungsbegriffe:
\begin{enumerate}[label=\Roman*.]
\item Partielle Ableitungen $\partial f/\partial x$ und $\partial f/\partial y$, wie im Reellen.
\item Komplexe Differenzierbarkeit:
\[
\lim_{z\to z_0}\frac{f(z)-f(z_0)}{z-z_0}\quad\text{existiert}.
\]
\item Holomorphie: komplex differenzierbar in $z_0$ und in einer offenen Menge um $z_0$.
\end{enumerate}

Setzt man $f=u+\ii v$ ein, erhält man
\[
f'(z_0)=\frac{\partial f}{\partial x}(z_0)=-\ii\frac{\partial f}{\partial y}(z_0).
\]

\begin{importantbox}
Die Cauchy-Riemann-Gleichungen lauten
\[
u_x(x,y)=v_y(x,y),\qquad u_y(x,y)=-v_x(x,y).
\]
Sind die partiellen Ableitungen in einer offenen Umgebung stetig und gelten dort die Cauchy-Riemann-Gleichungen, dann ist $f$ auf dieser Umgebung holomorph.
\end{importantbox}

\subsection{1.2 Beispiel zu Cauchy-Riemann}
\begin{beispielbox}
\[
f(z)=3x^3-9xy^2+\ii(9x^2y-3y^3).
\]
Dann
\[
u=3x^3-9xy^2,
\qquad
v=9x^2y-3y^3.
\]
Die partiellen Ableitungen sind
\[
u_x=9x^2-9y^2,
\quad u_y=-18xy,
\quad v_x=18xy,
\quad v_y=9x^2-9y^2.
\]
Also sind die Cauchy-Riemann-Gleichungen erfüllt und alle Ableitungen stetig. Daher ist $f$ holomorph.
\end{beispielbox}

\subsection{1.3 Beispiel: $f(z)=(x-\ii y)^2$}
\begin{fragebox}
Ist $f(z)=(x-\ii y)^2$ holomorph?
\end{fragebox}
\begin{antwortbox}
\[
f(z)=x^2-y^2-2\ii xy.
\]
Also
\[
u=x^2-y^2,
\qquad v=-2xy.
\]
Damit
\[
u_x=2x,
\quad u_y=-2y,
\quad v_x=-2y,
\quad v_y=-2x.
\]
Die Cauchy-Riemann-Gleichungen sind nur bei $(x,y)=(0,0)$ erfüllt. Dort ist $f$ komplex differenzierbar, aber nicht in einer offenen Umgebung. Also ist $f$ nirgends holomorph.
\end{antwortbox}

\begin{importantbox}
Falls eine Funktion wirklich von $\overline z$ abhängt, ist sie typischerweise nicht holomorph.
\end{importantbox}

\subsection{1.4 Harmonizität}
Sei $f=u+\ii v$ holomorph. Dann erfüllen $u$ und $v$ die Laplace-Gleichung:
\begin{importantbox}
\[
\frac{\partial^2u}{\partial x^2}+\frac{\partial^2u}{\partial y^2}=0,
\qquad
\frac{\partial^2v}{\partial x^2}+\frac{\partial^2v}{\partial y^2}=0.
\]
\end{importantbox}

\subsection{1.5 Bedingung an $c$}
\begin{fragebox}
Sei $f=u+\ii v$ und
\[
u=x^2+bxy+cy^2,
\qquad b,c\in\R.
\]
Unter welcher Bedingung ist $f$ holomorph?
\end{fragebox}
\begin{antwortbox}
Falls $f$ holomorph ist, ist $u$ harmonisch:
\[
u_{xx}+u_{yy}=2+2c=0.
\]
Also
\[
c=-1.
\]
\end{antwortbox}

\subsection{1.6 Holomorphie der konjugierten Funktion}
\begin{fragebox}
$f=u+\ii v$ ist holomorph. Ist $g(z)=\overline{f(\overline z)}$ holomorph?
\end{fragebox}
\begin{antwortbox}
Schreibt man $g(z)$ in Komponenten, erhält man
\[
g(x,y)=u(x,-y)-\ii v(x,-y).
\]
Mit den Cauchy-Riemann-Gleichungen für $f$ folgt, dass auch $g$ die Cauchy-Riemann-Gleichungen erfüllt. Auf dem entsprechend gespiegelten Definitionsgebiet ist $g$ daher holomorph.
\end{antwortbox}

\subsection{1.7 Komplexe Pfade}
Ein Pfad ist eine Abbildung
\[
\gamma:[a,b]\to\C.
\]
Man kann also Kurven in der komplexen Ebene parametrisieren.

\begin{center}
\begin{tikzpicture}[scale=1]
  \draw[axis] (-0.5,0)--(4,0) node[right] {$\Re$};
  \draw[axis] (0,-0.5)--(0,2.2) node[above] {$\Im$};
  \draw[blue,thick,patharrow] (0.5,0.5) .. controls (1.4,1.6) and (2.5,0.2) .. (3.3,1.5);
  \node at (0.45,0.25) {$\gamma(a)$};
  \node at (3.4,1.7) {$\gamma(b)$};
\end{tikzpicture}
\end{center}

\begin{beispielbox}
\[
\gamma:[0,1]\to\C,
\qquad
\gamma(t)=(1+\ii)t.
\]
\end{beispielbox}

\subsection{1.8 Kurvenintegral}
\begin{importantbox}
Für ein Kurvenintegral nehmen wir den Pfad $\gamma:[a,b]\to\C$ und definieren
\[
\int_\gamma f(z)\dd z=\int_a^b f(\gamma(t))\gamma'(t)\dd t.
\]
\end{importantbox}
Anschaulich wird der Vektor $f(\gamma(t))\gamma'(t)$ entlang des Pfades aufsummiert.

\subsection{2.0 Beispiel Kurvenintegral}
\begin{beispielbox}
Berechne
\[
\int_\gamma 3z\dd z,
\qquad \gamma(t)=1+\ii t,
\quad t\in[0,1].
\]
Dann $\gamma'(t)=\ii$ und
\[
\int_0^1 3(1+\ii t)\ii\dd t=\int_0^1(3\ii-3t)\dd t=3\ii-\frac32.
\]
\end{beispielbox}

\subsection{2.1 Parametrisierung einer Strecke}
\begin{fragebox}
Wie würde man den folgenden Weg parametrisieren: eine Strecke von $1$ nach $2$ auf der reellen Achse?
\end{fragebox}
\begin{antwortbox}
\[
\gamma:[0,1]\to\C,
\qquad
\gamma(t)=1+t.
\]
Allgemein gilt für eine gerade Strecke von $a$ nach $b$:
\begin{importantbox}
\[
\gamma(t)=a+t(b-a),\qquad t\in[0,1].
\]
\end{importantbox}
\end{antwortbox}

% =========================================================
\section{Woche 6: Kurvenintegrale}

\subsection{1.1 Parametrisierungen}
Für eine Strecke von $a$ nach $b$ verwendet man die Parametrisierung:
\[
\gamma_1(t)=a+t(b-a),\qquad t\in[0,1].
\]
Beispiel:
\[
\gamma_1(t)=-4+2\ii+t((-2+2\ii)-(-4+2\ii))=-4+2\ii+2t.
\]

Ein Kreis um $z_0$ mit Radius $r$ wird parametrisiert durch
\begin{importantbox}
\[
\gamma(t)=z_0+r\ee^{2\pi\ii t},\qquad t\in[0,1],
\]
oder alternativ durch $\gamma(t)=z_0+r\ee^{\ii t}$, $t\in[0,2\pi]$.
\end{importantbox}

\begin{center}
\begin{tikzpicture}[scale=0.8]
  \draw[axis] (-5,0)--(2,0) node[right] {$\Re$};
  \draw[axis] (0,-2)--(0,3) node[above] {$\Im$};
  \draw[blue,thick,->] (-4,2)--(-2,2) node[midway,above] {$\gamma_1$};
  \draw[purple,thick,patharrow] (0,0) circle (1);
  \node[purple] at (1.3,-0.4) {$\gamma_2$};
\end{tikzpicture}
\end{center}

\subsection{1.2 Parametrisierung einer geschlossenen Kurve}
\begin{fragebox}
Wie würde man die geschlossene obere Halbkreis-Kurve von $1$ nach $3$ mit Mittelpunkt $2$ parametrisieren?
\end{fragebox}
\begin{antwortbox}
\[
\gamma(t)=2-\ee^{-\pi\ii t},\qquad t\in[0,1].
\]
Für $t=0$ ist $\gamma(0)=1$, für $t=1$ ist $\gamma(1)=3$, und die Kurve läuft über den oberen Halbkreis.
\end{antwortbox}

\begin{center}
\begin{tikzpicture}[scale=0.9]
  \draw[axis] (0,0)--(4.5,0) node[right] {$\Re$};
  \draw[axis] (0,-1.2)--(0,2.2) node[above] {$\Im$};
  \draw[blue,thick,patharrow] (1,0) arc (180:0:1);
  \fill (1,0) circle (1.5pt) node[below] {$1$};
  \fill (2,0) circle (1pt) node[below] {$2$};
  \fill (3,0) circle (1.5pt) node[below] {$3$};
\end{tikzpicture}
\end{center}

\subsection{1.3 Wegunabhängigkeit und Stammfunktion}
Meist berechnet man geschlossene Wegintegrale:
\[
\int_\gamma f(z)\dd z=\int_0^1 f(\gamma(t))\gamma'(t)\dd t.
\]
Ist das Integral unabhängig vom genauen Weg und hängt nur von den Endpunkten ab, dann gilt
\[
\int_\gamma f(z)\dd z=F(\gamma(1))-F(\gamma(0)).
\]
Für geschlossene Wege ist dann das Integral gleich null.

\begin{importantbox}
Äquivalent dazu ist: $F'(z)=f(z)$, also besitzt $f$ die Stammfunktion $F$.
\end{importantbox}

Achtung: Das gilt nur, falls $f$ auf einer wegzusammenhängenden Menge definiert und stetig ist.

\subsection{1.4 Wegzusammenhängend}
\begin{importantbox}
Eine Menge heißt wegzusammenhängend, wenn es von jedem Punkt der Menge einen Pfad zu jedem anderen Punkt der Menge gibt.
\end{importantbox}

\begin{center}
\begin{tikzpicture}[scale=0.9]
  \draw[thick] (0,0) circle (0.8);
  \draw[thick] (1.8,0.1) circle (0.5);
  \node at (0,-1.1) {ja};
  \draw[thick] (4,0) circle (0.7);
  \draw[thick] (5.5,0) circle (0.7);
  \node at (4.75,-1.1) {nein};
  \draw[thick] (8,0) .. controls (8.7,0.8) and (9.5,0.5) .. (9.3,-0.4) .. controls (8.8,-1) and (7.6,-0.8) .. cycle;
  \node at (8.7,-1.1) {ja};
\end{tikzpicture}
\end{center}

\subsection{1.5 Einfach zusammenhängend}
\begin{importantbox}
Eine Menge ist einfach zusammenhängend, wenn sie wegzusammenhängend ist und man jede geschlossene Schlaufe in der Menge auf einen Punkt zusammenziehen kann.
\end{importantbox}

Beispiele:
\begin{enumerate}[label=\alph*)]
\item $\R^2\setminus\{0\}$: nicht einfach zusammenhängend.
\item $\{z\in\C\mid \Imm(z)\neq0\}$: nicht wegzusammenhängend.
\item $\{z\in\C\mid |\Imm(z)|>0\text{ und }|z|>3\}$: je nach Gebietskomponente zu betrachten.
\item $\{z\in\C\mid \Imm(z)>0\text{ oder }\Imm(z)+\Ree(z)>0\}$: einfach zusammenhängend.
\item $\{z\in\C\mid 0<|z|<1\}$: nicht einfach zusammenhängend.
\item $\{z\in\C\mid 2<|z|<3\}$: nicht einfach zusammenhängend.
\end{enumerate}

\subsection{1.6 Kurvenintegral über $z^2$}
\begin{fragebox}
Was ist
\[
\int_\gamma z^2\dd z,
\qquad \gamma(t)=\ee^{\pi\ii t},\quad t\in[0,1]?
\]
\end{fragebox}
\begin{antwortbox}
Da $F(z)=\frac13z^3$ eine Stammfunktion von $z^2$ ist,
\[
\int_\gamma z^2\dd z=F(\gamma(1))-F(\gamma(0))
=F(-1)-F(1)=-\frac13-\frac13=-\frac23.
\]
\end{antwortbox}

\subsection{1.7 Betrag im Kurvenintegral}
\begin{fragebox}
Was ist
\[
\int_\gamma |z|^2\dd z
\]
für den Weg, der von $-1$ nach $1$ über die reelle Strecke und dann über den oberen Halbkreis zurückläuft?
\end{fragebox}
\begin{antwortbox}
Der Weg wird in zwei Teile zerlegt.

Erster Weg: $\gamma_1(t)=-1+2t$, $t\in[0,1]$. Dann
\[
I_1=\int_0^1|-1+2t|^2\cdot2\dd t=\frac23.
\]
Zweiter Weg: $\gamma_2(t)=\ee^{\pi\ii t}$, $t\in[0,1]$. Dann $|\gamma_2(t)|^2=1$ und
\[
I_2=\int_0^1 \pi\ii\ee^{\pi\ii t}\dd t=\ee^{\pi\ii}-1=-2.
\]
Also
\[
\int_\gamma |z|^2\dd z=\frac23-2=-\frac43.
\]
Man muss $|z|^2$ nicht als $z\overline z$ behandeln, wenn man direkt parametrisiert.
\end{antwortbox}

\subsection{1.8 Cauchyscher Integralsatz}
\begin{importantbox}
Falls $U$ einfach zusammenhängend ist und $f:U\to\C$ holomorph, dann besitzt $f$ eine Stammfunktion. Für jeden geschlossenen Weg $\gamma$ gilt
\[
\int_\gamma f(z)\dd z=0.
\]
\end{importantbox}

% =========================================================
\section{Woche 7: Cauchysche Integralformel}

\subsection{1.1 Cauchysche Integralformel}
Direkte Parametrisierung von Kurvenintegralen kann schnell kompliziert werden. Deshalb verwendet man die Cauchysche Integralformel.

\begin{importantbox}
Ist $U$ einfach zusammenhängend, $\gamma$ geschlossen, positiv orientiert, $z_0\in U$ und $f$ holomorph, dann gilt
\[
f(z_0)=\frac{1}{2\pi\ii}\int_\gamma \frac{f(z)}{z-z_0}\dd z.
\]
\end{importantbox}

\begin{center}
\begin{tikzpicture}[scale=1]
  \draw[axis] (-1.2,0)--(4,0) node[right] {$\Re$};
  \draw[axis] (0,-1.2)--(0,3) node[above] {$\Im$};
  \draw[purple,dashed] (2,1.1) circle (1.3);
  \draw[blue,thick,patharrow] (2,1.1) +(-0.8,0) .. controls (1.3,2) and (2.4,2.5) .. (2.9,1.6) .. controls (3.2,0.7) and (1.8,0.4) .. cycle;
  \fill (2,1.1) circle (1.5pt) node[right] {$z_0$};
  \node[blue] at (3.2,1.7) {$\gamma$ läuft gegen den Uhrzeigersinn};
  \node[purple] at (3.4,2.4) {$U$};
\end{tikzpicture}
\end{center}

\subsection{1.2 Beispiel}
\begin{beispielbox}
Berechne
\[
\int_\gamma \frac{1}{z(z+3\ii)}\dd z,
\qquad \gamma(t)=\ee^{2\pi\ii t}.
\]
\end{beispielbox}
\begin{kochrezeptbox}
\begin{enumerate}[label=\Roman*.]
\item Bestimme $z_0$: Im Nenner steht $z$, also $z_0=0$. Außerdem liegt $-3\ii$ nicht in $\gamma$.
\item Schreibe
\[
\frac{1}{z(z+3\ii)}=\frac{f(z)}{z-z_0},\qquad f(z)=\frac1{z+3\ii}.
\]
\item Dann
\[
\int_\gamma \frac{f(z)}{z}\dd z=2\pi\ii f(0)=2\pi\ii\frac1{3\ii}=\frac{2\pi}{3}.
\]
\end{enumerate}
\end{kochrezeptbox}

\subsection{1.3 Orientierung}
\begin{fragebox}
Wie verändert sich das Integral, wenn $\gamma$ im Uhrzeigersinn läuft?
\end{fragebox}
\begin{antwortbox}
Das Vorzeichen ändert sich. Man multipliziert das Ergebnis mit $-1$.
\[
\int_{-\gamma} f(z)\dd z=-\int_\gamma f(z)\dd z.
\]
\end{antwortbox}

\subsection{1.4 Allgemeine Cauchy-Formel für Ableitungen}
\begin{importantbox}
Für $n\in\N_0$ gilt
\[
f^{(n)}(z_0)=\frac{n!}{2\pi\ii}\int_\gamma\frac{f(z)}{(z-z_0)^{n+1}}\dd z.
\]
Das ist eine der wichtigsten Formeln der Vorlesung. Es folgt auch: Ist $f$ holomorph, dann ist $f$ unendlich oft komplex differenzierbar.
\end{importantbox}

\subsection{1.5 Beispiel mit Ableitungsformel}
\begin{beispielbox}
Berechne
\[
\int_\gamma \frac{\ee^{3z}}{(z-1)^2}\dd z,
\qquad \gamma(t)=3\ee^{2\pi\ii t}.
\]
\end{beispielbox}
Hier ist $z_0=1$ und $n=1$, $f(z)=\ee^{3z}$. Daher
\[
\int_\gamma \frac{f(z)}{(z-1)^2}\dd z=\frac{2\pi\ii}{1!}f'(1)
=2\pi\ii\cdot3\ee^3=6\pi\ii\ee^3.
\]

\subsection{1.6 Frage mit $\ee^{-2/z}$}
\begin{fragebox}
Was ist
\[
\int_\gamma \frac{\ee^{-2/z}}{(z+3\ii)^2}\dd z,
\qquad \gamma(t)=4\ee^{2\pi\ii t}?
\]
\end{fragebox}
\begin{antwortbox}
Hier ist $z_0=-3\ii$ und $n=1$. Setze $f(z)=\ee^{-2/z}$. Dann
\[
\int_\gamma \frac{f(z)}{(z+3\ii)^2}\dd z=2\pi\ii f'(-3\ii).
\]
Da
\[
f'(z)=\ee^{-2/z}\cdot\frac{2}{z^2},
\]
folgt
\[
2\pi\ii\cdot \ee^{-2/(-3\ii)}\cdot\frac{2}{(-3\ii)^2}.
\]
\end{antwortbox}

\subsection{1.7 Mittelwertsatz aus der Cauchy-Formel}
Für einen Kreis $\gamma(t)=z_0+R\ee^{2\pi\ii t}$ ergibt die Cauchy-Formel den Mittelwertsatz:
\begin{importantbox}
\[
f(z_0)=\int_0^1 f\bigl(z_0+R\ee^{2\pi\ii t}\bigr)\dd t.
\]
\end{importantbox}
Dieser Satz wird häufig für Beweise verwendet.

\subsection{1.8 Maximum-Modulus-Prinzip}
Aus dem Mittelwertsatz folgt:
\begin{importantbox}
Ist $f$ holomorph auf $U$, $U$ offen und wegzusammenhängend, dann hat $|f|$ kein Maximum im Inneren von $U$, außer $f$ ist konstant. Das Maximum von $|f|$ liegt auf dem Rand von $U$.
\end{importantbox}

\begin{center}
\begin{tikzpicture}[scale=0.9]
  \draw[axis] (-1,0)--(5,0) node[right] {$\Re$};
  \draw[axis] (0,-1)--(0,4) node[above] {$\Im$};
  \filldraw[blue!20,draw=blue!70!black,thick] (2,2) .. controls (1.3,3) and (3.3,3.6) .. (3.8,2.5) .. controls (4.2,1.2) and (2.1,0.8) .. (1.3,1.5) .. controls (0.8,2) and (1.3,2.7) .. cycle;
  \fill (3.5,2.3) circle (2pt) node[right] {$\max |f(z)|$ auf dem Rand};
  \node at (2.5,3.5) {$U$};
\end{tikzpicture}
\end{center}

\begin{beispielbox}
Sei $U=B(0,1)$ und $f(z)=\cos z$. Dann liegt das Maximum von $|\cos z|$ auf dem Einheitskreis.
\end{beispielbox}

% =========================================================
\section{Woche 8: Reihen}

\subsection{1.1 Taylorreihe}
Die reelle Taylorreihe ist bekannt:
\[
f(x)=\sum_{n=0}^{\infty}\frac{f^{(n)}(x_0)}{n!}(x-x_0)^n,
\]
falls $f$ unendlich oft differenzierbar ist.

Im Komplexen ist es ähnlich:
\begin{importantbox}
Ist $f$ holomorph, dann gilt
\[
f(z)=\sum_{n=0}^{\infty}\frac{f^{(n)}(z_0)}{n!}(z-z_0)^n.
\]
Holomorphie reicht, da holomorph impliziert: unendlich oft komplex differenzierbar.
\end{importantbox}

\subsection{1.2 Geometrische Reihe}
Wir benutzen häufig
\begin{importantbox}
\[
\frac1{1-z}=\sum_{n=0}^{\infty}z^n,
\qquad |z|<1.
\]
\end{importantbox}
Der Gedanke ist: Statt $z$ darf man auch einen Ausdruck $f(z)$ einsetzen:
\[
\frac1{1-f(z)}=\sum_{n=0}^{\infty}(f(z))^n,
\qquad |f(z)|<1.
\]

\subsection{1.3 Reihenentwicklung durch Umformen}
Praktisch entwickelt man viele Funktionen, indem man sie so umformt, dass die geometrische Reihe direkt anwendbar ist.

\begin{beispielbox}
Entwickle
\[
\frac1{1+z}
\]
um den Punkt $z_0=0$.
\end{beispielbox}
\[
\frac1{1+z}=\frac1{1-(-z)}=\sum_{n=0}^{\infty}(-z)^n
=\sum_{n=0}^{\infty}(-1)^nz^n,
\qquad |z|<1.
\]
Der Ausdruck $(z-z_0)^n$ zeigt, um welchen Punkt entwickelt wird.

Nun um $z_0=1$:
\[
\frac1{1+z}=\frac1{1+(z-1)+1}=\frac1{2+(z-1)}
=\frac12\frac1{1+\frac{z-1}{2}}.
\]
Daher
\[
\frac1{1+z}=\frac12\sum_{n=0}^{\infty}\left(-\frac{z-1}{2}\right)^n
=\sum_{n=0}^{\infty}\frac{(-1)^n}{2^{n+1}}(z-1)^n,
\]
mit Konvergenzbedingung $|z-1|<2$.

\begin{importantbox}
Allgemein:
\[
\frac1{z+a}=\frac1{a+z_0}\cdot\frac1{1+\frac{z-z_0}{a+z_0}},
\]
und dann wird die geometrische Reihe verwendet.
\end{importantbox}

% =========================================================
\section{Woche 9: Reihen und Residuum}

\subsection{1.1 Taylorkoeffizienten und Problem}
Für Taylorreihen gilt:
\[
f(z)=\sum_{n=0}^{\infty}\frac{f^{(n)}(z_0)}{n!}(z-z_0)^n
=\sum_{n=0}^{\infty}a_n(z-z_0)^n.
\]
Außerdem gilt nach der Cauchy-Formel:
\[
f^{(n)}(z_0)=\frac{n!}{2\pi\ii}\int_\gamma\frac{f(z)}{(z-z_0)^{n+1}}\dd z.
\]
Daher
\begin{importantbox}
\[
a_n=\frac{1}{n!}f^{(n)}(z_0)
=\frac1{2\pi\ii}\int_\gamma\frac{f(z)}{(z-z_0)^{n+1}}\dd z.
\]
\end{importantbox}
Problem: Diese Darstellung gilt nur auf einer Kreisscheibe, auf der $f$ holomorph ist.

\subsection{1.2 Laurent-Reihe}
Die Lösung ist, auch negative Potenzen zuzulassen. Dann gilt die Reihe auf einem Ring.

\begin{importantbox}
Die Laurent-Reihe um $z_0$ lautet
\[
f(z)=\sum_{n=-\infty}^{\infty}c_n(z-z_0)^n,
\qquad
c_n=\frac1{2\pi\ii}\int_\gamma\frac{f(z)}{(z-z_0)^{n+1}}\dd z.
\]
\end{importantbox}
Wenn wir $c_n$ kennen, kennen wir auch Integrale. Insbesondere
\begin{importantbox}
\[
\int_\gamma f(z)\dd z=2\pi\ii\,c_{-1}.
\]
\end{importantbox}

\begin{center}
\begin{tikzpicture}[scale=0.95]
  \draw[axis] (-2.5,0)--(2.5,0) node[right] {$\Re$};
  \draw[axis] (0,-2.5)--(0,2.5) node[above] {$\Im$};
  \fill[red!15] (0,0) circle (1.8);
  \fill[white] (0,0) circle (0.8);
  \draw[red!60!black,thick] (0,0) circle (1.8);
  \draw[red!60!black,thick] (0,0) circle (0.8);
  \node at (1.35,1.35) {Ring};
\end{tikzpicture}
\end{center}

\subsection{1.3 Beispiel: Integral von $\ee^{1/z}$}
\begin{beispielbox}
Berechne
\[
\int_\gamma \ee^{1/z}\dd z,
\qquad \gamma(t)=\ee^{2\pi\ii t}.
\]
\end{beispielbox}
\[
\ee^{1/z}=\sum_{n=0}^{\infty}\frac{1}{n!z^n}
=1+\frac1z+\frac1{2z^2}+\cdots.
\]
Der Koeffizient $c_{-1}$ ist $1$. Also
\[
\int_\gamma \ee^{1/z}\dd z=2\pi\ii.
\]
Bemerkung: $\ee^{1/z}$ ist auf dem Ring $\C\setminus\{0\}$ definiert; daher wird die Laurent-Reihe verwendet.

\subsection{1.4 Frage}
\begin{fragebox}
Was ergibt
\[
\int_\gamma \ee^{1/z^2}\dd z?
\]
\end{fragebox}
\begin{antwortbox}
\[
\ee^{1/z^2}=\sum_{n=0}^{\infty}\frac1{n!z^{2n}}
=1+\frac1{z^2}+\frac1{2z^4}+\cdots.
\]
Es gibt keinen $z^{-1}$-Term, also ist $c_{-1}=0$ und
\[
\int_\gamma \ee^{1/z^2}\dd z=0.
\]
\end{antwortbox}

\subsection{1.5 Residuum}
\begin{importantbox}
Der Koeffizient $c_{-1}$ ist so besonders, dass er einen eigenen Namen hat:
\[
c_{-1}=\Res(f,z_0).
\]
\end{importantbox}

\subsection{1.6 Laurent-Reihe von $\ee^{1/z}/z$}
\begin{fragebox}
Was ist die Laurent-Reihe von
\[
f(z)=\frac{\ee^{1/z}}{z}?
\]
Was ist
\[
\int_\gamma f(z)\dd z,
\qquad \gamma(t)=\ee^{2\pi\ii t}?
\]
\end{fragebox}
\begin{antwortbox}
\[
\frac{\ee^{1/z}}{z}=\frac1z\sum_{n=0}^{\infty}\frac1{n!z^n}
=\sum_{n=0}^{\infty}\frac1{n!}z^{-(n+1)}.
\]
Der Koeffizient von $z^{-1}$ ist $1$, also
\[
\int_\gamma f(z)\dd z=2\pi\ii.
\]
\end{antwortbox}

\subsection{1.7 Arten von Singularitäten}
Sei
\[
f(z)=\sum_{n=-\infty}^{\infty}c_n(z-z_0)^n.
\]
Dann gibt es drei Arten von isolierten Singularitäten:
\begin{enumerate}[label=\Roman*.]
\item Wesentlich: $c_n\neq0$ für unendlich viele negative $n$; Beispiel $\sin(1/z)$.
\item Pol der Ordnung $m$: $c_{-m}\neq0$ und $c_n=0$ für alle $n<-m$; Beispiel $1/z^3$ hat Polordnung $3$ bei $0$.
\item Hebbar: $c_n=0$ für alle $n<0$; dann kann die Funktion fortgesetzt werden.
\end{enumerate}

\subsection{1.8 Singularitäten bestimmen}
\begin{fragebox}
Bestimme die Singularitäten und deren Art von
\[
f(z)=\frac{\cos(3/z)\sin(z-1)}{(z-4)(z^2+1)(z-1)}.
\]
\end{fragebox}
\begin{kochrezeptbox}
\begin{enumerate}[label=\Roman*.]
\item Alle möglichen Problemstellen aus Nenner und Funktionsausdrücken sammeln.
\item Prüfen, ob sich Faktoren kürzen oder Grenzwerte endlich sind.
\item Nach wesentlicher Singularität, Pol oder hebbarer Singularität klassifizieren.
\end{enumerate}
\end{kochrezeptbox}
\begin{antwortbox}
Die möglichen Problemstellen kommen aus dem Nenner und aus dem Term $\cos(3/z)$:
\[
z=0,\qquad z=4,\qquad z=\ii,\qquad z=-\ii,\qquad z=1.
\]
Bei $z=0$ liegt wegen $\cos(3/z)$ eine wesentliche Singularität vor. Bei $z=4$, $z=\ii$ und $z=-\ii$ liegen einfache Pole vor, da dort jeweils eine einfache Nullstelle des Nenners steht und der Zähler im Allgemeinen nicht verschwindet. Bei $z=1$ heben sich $\sin(z-1)$ und $(z-1)$ gegenseitig auf, denn
\[
\lim_{z\to1}\frac{\sin(z-1)}{z-1}=1.
\]
Daher ist $z=1$ eine hebbare Singularität.
\end{antwortbox}

\subsection{1.9 Laurent-Reihe auf einem Ring}
Die Laurent-Reihe gilt nur auf einem Ring. Will man $f$ auf unterschiedlichen Bereichen darstellen, braucht man separate Reihen pro Teilgebiet.

\begin{beispielbox}
Für
\[
f(z)=\frac1{(z-1)(z-4)}
\]
werden die Bereiche
\[
I: |z|<1,
\qquad II: 1<|z|<4,
\qquad III: |z|>4
\]
betrachtet.
\end{beispielbox}

\begin{center}
\begin{tikzpicture}[scale=0.85]
  \draw[axis] (-1.5,0)--(5.3,0) node[right] {$\Re$};
  \draw[axis] (0,-2.3)--(0,2.3) node[above] {$\Im$};
  \draw[blue] (0,0) circle (1);
  \draw[blue] (0,0) circle (4);
  \fill (1,0) circle (2pt) node[below] {$1$};
  \fill (4,0) circle (2pt) node[below] {$4$};
  \node at (0.4,0.4) {I};
  \node at (2.4,1.2) {II};
  \node at (4.5,1.8) {III};
\end{tikzpicture}
\end{center}

\subsection{2.1 Laurent-Reihe im Gebiet I}
\begin{fragebox}
Bestimme die Laurent-Reihe von $f$ im Gebiet I, also $|z|<1$.
\end{fragebox}
\begin{antwortbox}
Partialbruchzerlegung:
\[
\frac1{(z-1)(z-4)}=\frac{A}{z-1}+\frac{B}{z-4}
=-\frac1{3}\frac1{z-1}+\frac1{3}\frac1{z-4}.
\]
Für $|z|<1$:
\[
\frac1{z-1}=-\frac1{1-z}=-\sum_{n=0}^{\infty}z^n,
\]
\[
\frac1{z-4}=-\frac14\frac1{1-z/4}=-\frac14\sum_{n=0}^{\infty}\left(\frac z4\right)^n.
\]
Zusammen ergibt sich eine Taylor-/Laurent-Reihe um $0$ im Bereich $|z|<1$.
\end{antwortbox}

% =========================================================
\section{Woche 10: Residuensatz}

\subsection{1.1 Mehrere Singularitäten}
Wenn $f$ mehrere isolierte Singularitäten innerhalb einer geschlossenen Kurve besitzt, summiert man die Beiträge.

\begin{importantbox}
Residuensatz:
\[
\int_\gamma f(z)\dd z=2\pi\ii\sum_{j=1}^{n}\Res(f,z_j),
\]
wobei die Singularitäten $z_1,\ldots,z_n$ von $\gamma$ umschlossen werden.
\end{importantbox}

\begin{center}
\begin{tikzpicture}[scale=0.9]
  \draw[axis] (-0.5,0)--(5.2,0) node[right] {$\Re$};
  \draw[axis] (0,-1.5)--(0,2.5) node[above] {$\Im$};
  \draw[blue,thick,patharrow] (1,0.5) .. controls (1.2,2) and (3.4,2.3) .. (4.2,1.1) .. controls (5,-0.3) and (2.7,-1) .. (1.2,-0.5) .. controls (0.4,-0.1) and (0.5,0.7) .. cycle;
  \draw[purple,patharrow] (1.7,0.5) circle (0.45);
  \draw[purple,patharrow] (2.6,0.7) circle (0.45);
  \fill (1.7,0.5) circle (2pt) node[left] {$z_0$};
  \fill (2.6,0.7) circle (2pt) node[right] {$z_1$};
\end{tikzpicture}
\end{center}

\subsection{1.2 Windungszahl}
Genauer gilt:
\begin{importantbox}
\[
\int_\gamma f(z)\dd z=2\pi\ii\sum_j W(\gamma,z_j)\Res(f,z_j).
\]
\end{importantbox}
Die Windungszahl $W(\gamma,z_j)$ sagt, ob $z_j$ positiv, also gegen den Uhrzeigersinn, oder negativ, also im Uhrzeigersinn, und wie oft umschlossen wird.

\subsection{1.3 Beispiel Windungszahl}
Eine Kurve, die einen Punkt zweimal im Uhrzeigersinn umläuft, hat Windungszahl $-2$.

\subsection{1.4 Windungszahlen bestimmen}
\begin{fragebox}
Bestimme die Windungszahlen für die gezeichneten Kurven und Punkte.
\end{fragebox}
\begin{antwortbox}
Die Ergebnisse lauten:
\begin{enumerate}[label=\alph*)]
\item Ein Punkt wird einmal im Uhrzeigersinn umlaufen: $W(\gamma,z_0)=-1$.
\item Zwei Punkte: $W(\gamma,z_0)=0$, $W(\gamma,z_1)=-1$.
\item Ein Punkt wird zweimal gegen den Uhrzeigersinn umlaufen: $W(\gamma,z_0)=2$.
\end{enumerate}
\end{antwortbox}

\subsection{1.5 Beispiel mit Residuensatz}
\begin{beispielbox}
Bestimme
\[
\int_\gamma \frac{1}{z^2-2z}\dd z
\]
für die gezeichnete komplizierte Kurve.
\end{beispielbox}
\begin{kochrezeptbox}
\begin{enumerate}[label=\Roman*.]
\item Bestimmen, welche Singularitäten innerhalb von $\gamma$ liegen und welche Ordnung sie haben.
\item Windungszahlen der Kurve um alle relevanten Singularitäten bestimmen.
\item Für jede Singularität das Residuum berechnen.
\item Residuensatz mit Windungszahlen anwenden.
\end{enumerate}
\end{kochrezeptbox}
Die Singularitäten sind
\[
\frac1{z(z-2)},
\]
also bei $z_0=0$ und $z_1=2$, beide von Ordnung $1$.
Die Windungszahlen aus der Zeichnung sind
\[
W(\gamma,0)=2,
\qquad W(\gamma,2)=1.
\]
Für einfache Pole gilt:
\begin{importantbox}
Wenn
\[
f(z)=\frac{\varphi(z)}{(z-z_j)^m},
\]
wobei $\varphi$ holomorph bei $z_j$ ist, dann
\[
\Res(f,z_j)=\frac{\varphi^{(m-1)}(z_j)}{(m-1)!}.
\]
Für $m=1$ ist $\Res(f,z_j)=\varphi(z_j)$.
\end{importantbox}

Hier:
\[
\Res(f,0)=\lim_{z\to0}\frac{z}{z(z-2)}=-\frac12,
\]
\[
\Res(f,2)=\lim_{z\to2}\frac{z-2}{z(z-2)}=\frac12.
\]
Also
\[
\int_\gamma f(z)\dd z
=2\pi\ii\left(2\cdot\left(-\frac12\right)+1\cdot\frac12\right)
=-\pi\ii.
\]

\subsection{1.6 Beispiel mit Laurent-Reihe}
\begin{beispielbox}
Berechne
\[
\int_\gamma \frac{\sin(1/z)}{z^2}\dd z.
\]
\end{beispielbox}
Hier erkennt man nicht sofort einen Pol endlicher Ordnung bei $0$. Man entwickelt:
\[
\sin(1/z)=\frac1z-\frac1{3!z^3}+\frac1{5!z^5}-\cdots.
\]
Dann
\[
\frac{\sin(1/z)}{z^2}=\frac1{z^3}-\frac1{3!z^5}+\frac1{5!z^7}-\cdots.
\]
Es gibt keinen $z^{-1}$-Term, also $\Res(f,0)=0$ und das Integral ist $0$.

\subsection{1.7 Residuenfrage}
\begin{fragebox}
Was ist $\Res(f,0)$ für
\[
f(z)=\frac{z^2\sin z}{z}\,?
\]
Was ist $\Res(f,2)$ für
\[
f(z)=\frac{\ee^z}{z-2}?
\]
\end{fragebox}
\begin{tippbox}
Führe die Schritte I und III des Kochrezepts aus: Art der Singularität erkennen und dann das Residuum berechnen.
\end{tippbox}
\begin{antwortbox}
Für das erste Beispiel wird per Laurent-/Taylorentwicklung geprüft, ob ein $z^{-1}$-Term auftritt. Da $z\sin z$ holomorph ist, ist $\Res(f,0)=0$.

Für das zweite Beispiel liegt ein einfacher Pol bei $z=2$ vor:
\[
\Res\left(\frac{\ee^z}{z-2},2\right)=\ee^2.
\]
\end{antwortbox}

\subsection{1.8 Identitätssatz}
\begin{importantbox}
Falls $f$ holomorph ist auf $U$ und $f(z)=0$ auf einer offenen Teilmenge oder auf einer Menge mit Häufungspunkt in $U$, dann gilt $f(z)=0$ überall auf $U$.
\end{importantbox}
Daraus folgt: Falls $f$ und $g$ holomorph sind und $f(z)=g(z)$ auf einer solchen Menge gilt, dann gilt $f=g$ auf $U$.

\subsection{1.9 Beispiel Identitätssatz}
Wir wissen, dass
\[
\sin^2 x+\cos^2 x=1
\]
für alle reellen $x$. Da beide Seiten holomorphe Funktionen sind, gilt die Identität auch im Komplexen:
\[
\sin^2 z+\cos^2 z=1.
\]

% =========================================================
\section{Woche 11: Anwendung auf reelle Integrale}

\subsection{Idee}
Die zentrale Idee ist:
\begin{center}
reelles Integral $\longrightarrow$ komplexer Raum $\longrightarrow$ Residuensatz $\longrightarrow$ zurück in den reellen Raum.
\end{center}

\subsection{1.1 Erster Fall: Integrale über $[0,2\pi]$}
Wir betrachten Integrale der Form
\[
\int_0^{2\pi}F(\cos t,\sin t)\dd t.
\]
Substituiere
\begin{importantbox}
\[
z=\ee^{\ii t},
\qquad
\dd z=\ii\ee^{\ii t}\dd t=\ii z\dd t,
\qquad
\cos t=\frac{z+z^{-1}}{2},
\qquad
\sin t=\frac{z-z^{-1}}{2\ii}.
\]
\end{importantbox}

\begin{beispielbox}
Bestimme
\[
\int_0^{2\pi}\frac1{5-4\cos t}\dd t.
\]
\end{beispielbox}
\begin{kochrezeptbox}
\begin{enumerate}[label=\Roman*.]
\item Substituiere $z=\ee^{\ii t}$ und ersetze $\cos t$, $\sin t$ sowie $\dd t$.
\item Bestimme alle Singularitäten der entstehenden Funktion und prüfe, welche im Einheitskreis liegen.
\item Berechne die Residuen der Singularitäten im Einheitskreis.
\item Wende den Residuensatz an und vereinfache das Ergebnis.
\end{enumerate}
\end{kochrezeptbox}
Nach Substitution erhält man ein Kurvenintegral über den Einheitskreis. Die Singularitäten ergeben sich aus
\[
-2\ii z^2+5\ii z-2\ii=0
\quad\Longleftrightarrow\quad
z^2-\frac52z+1=0.
\]
Daraus
\[
z_1=\frac12,
\qquad z_2=2.
\]
Nur $z_1=1/2$ liegt im Einheitskreis. Der Residuensatz liefert
\[
\int_0^{2\pi}\frac1{5-4\cos t}\dd t=\frac{2\pi}{3}.
\]

\subsection{1.2 Frage}
\begin{fragebox}
Was ergibt
\[
\int_0^{2\pi}\frac{3}{5\sin t+3}\dd t?
\]
\end{fragebox}
\begin{antwortbox}
Zuerst muss geprüft werden, ob das reelle Integral überhaupt existiert. Da
\[
5\sin t+3=0
\quad\Longleftrightarrow\quad
\sin t=-\frac35
\]
für zwei Werte $t\in[0,2\pi]$ erfüllt ist, hat der Integrand im Integrationsintervall Polstellen. Das uneigentliche reelle Integral existiert daher nicht im gewöhnlichen Sinn.

Die formale Substitution $z=\ee^{\ii t}$ führt auf Pole bei
\[
z_1=\ii(-3+2\sqrt2),
\qquad
z_2=\ii(-3-2\sqrt2),
\]
wobei nur $z_1$ im Einheitskreis liegt. Diese Rechnung ist jedoch nur mit zusätzlicher Interpretation, etwa als Hauptwert mit geeigneter Konturbehandlung, sinnvoll. Für das gewöhnliche reelle Integral ist die korrekte Aussage: Es divergiert.
\end{antwortbox}

\subsection{1.3 Zweiter Fall: Integrale über $\R$}
Für Integrale
\[
\int_{-\infty}^{\infty}f(t)\dd t
\]
verwendet man den oberen Halbkreis $\gamma_R$.

\begin{center}
\begin{tikzpicture}[scale=0.85]
  \draw[axis] (-3.8,0)--(3.8,0) node[right] {$\Re$};
  \draw[axis] (0,-0.5)--(0,3.5) node[above] {$\Im$};
  \draw[blue,thick,patharrow] (-3,0) arc (180:0:3);
  \draw[green!60!black,thick,->] (-3,0)--(3,0) node[midway,below] {$\gamma_R^-$};
  \fill (1.2,1.5) circle (2pt) node[right] {$z_1$};
  \fill (-1.0,1.1) circle (2pt) node[left] {$z_2$};
  \node at (3,-0.25) {$R$};
  \node at (-3,-0.25) {$-R$};
\end{tikzpicture}
\end{center}

Wenn $f=p/q$ rational ist, $\deg q\ge \deg p+2$, $q$ keine reellen Nullstellen hat und $f$ beschränkt ist, dann verschwindet der Anteil über dem Halbkreis im Grenzwert $R\to\infty$ und
\begin{importantbox}
\[
\int_{-\infty}^{\infty}f(t)\dd t=2\pi\ii\sum_{\Imm z_j>0}\Res(f,z_j).
\]
\end{importantbox}

\subsection{1.4 Beispiel}
\begin{beispielbox}
Bestimme
\[
\int_{-\infty}^{\infty}\frac{\ee^{\ii t}}{t^2-6\ii t+13}\dd t.
\]
\end{beispielbox}
\begin{kochrezeptbox}
\begin{enumerate}[label=\Roman*.]
\item Prüfe, ob der Halbkreis-Trick anwendbar ist: keine reellen Pole und ausreichendes Abklingen.
\item Bestimme die Pole in der oberen Halbebene.
\item Berechne die Residuen an diesen Polen.
\item Wende den Residuensatz an. Der Halbkreisanteil verschwindet im Grenzwert.
\end{enumerate}
\end{kochrezeptbox}
Die Pole sind
\[
z_1=3+2\ii,
\qquad z_2=3-2\ii.
\]
Nur $z_1$ liegt in der oberen Halbebene. Es gilt
\[
\Res(f,z_1)=\lim_{z\to z_1}(z-z_1)\frac{\ee^{\ii z}}{(z-z_1)(z-z_2)}
=\frac{\ee^{\ii(3+2\ii)}}{(3+2\ii)-(3-2\ii)}
=\frac{\ee^{-2+3\ii}}{4\ii}.
\]
Daher
\[
\int_{-\infty}^{\infty}\frac{\ee^{\ii t}}{t^2-6\ii t+13}\dd t
=2\pi\ii\frac{\ee^{-2+3\ii}}{4\ii}
=\frac{\pi}{2}\ee^{-2}(\cos3+\ii\sin3).
\]

\subsection{1.5 Real- und Imaginärteil}
\begin{fragebox}
Was ergibt
\[
\int_{-\infty}^{\infty}\frac{\cos(5t)}{t^2-6\ii t+13}\dd t?
\]
Was ergibt
\[
\int_{-\infty}^{\infty}\frac{\sin(5t)}{t^2-6\ii t+13}\dd t?
\]
\end{fragebox}
\begin{antwortbox}
Man berechnet zunächst das Integral mit $\ee^{\ii 5t}$ und nimmt danach Real- bzw. Imaginärteil:
\[
\int_{-\infty}^{\infty}\frac{\cos(5t)}{t^2-6\ii t+13}\dd t
=\Ree\left(\int_{-\infty}^{\infty}\frac{\ee^{\ii5t}}{t^2-6\ii t+13}\dd t\right),
\]
\[
\int_{-\infty}^{\infty}\frac{\sin(5t)}{t^2-6\ii t+13}\dd t
=\Imm\left(\int_{-\infty}^{\infty}\frac{\ee^{\ii5t}}{t^2-6\ii t+13}\dd t\right).
\]
Aus der Rechnung ergibt sich entsprechend ein Faktor
\[
\frac{\pi}{2}\ee^{-10}\cos(15)
\quad\text{bzw.}\quad
\frac{\pi}{2}\ee^{-10}\sin(15).
\]
\end{antwortbox}

\subsection{1.6 Abschätzung des Halbkreisanteils}
\begin{fragebox}
Was ist die obere Abschätzung von
\[
\left|\int_{\gamma_R}\frac{\ee^{\ii z}}{z^2-6\ii z+13}\dd z\right|?
\]
\end{fragebox}
\begin{antwortbox}
Auf dem oberen Halbkreis gilt $|\ee^{\ii z}|\le1$. Außerdem ist
\[
|z^2-6\ii z+13|\ge R^2-6R+13.
\]
Die Länge des Halbkreises ist $\pi R$. Also
\[
\left|\int_{\gamma_R}\frac{\ee^{\ii z}}{z^2-6\ii z+13}\dd z\right|
\le \frac{\pi R}{R^2-6R+13}\xrightarrow{R\to\infty}0.
\]
\end{antwortbox}

\subsection{1.7 Anschauliche Einführung in Fourier-Analyse}
Differenzierbare Funktionen kann man als Taylorreihe schreiben:
\[
f(x)=\sum_{n=0}^{\infty}a_nx^n.
\]
Dabei sind $1,x,x^2,\ldots$ die Basisfunktionen.

Wenn $f$ periodisch ist, wählt man eine periodische Basis:
\[
f(t)=\frac{a_0}{2}+\sum_{n=1}^{\infty}a_n\cos\left(\frac{n\pi t}{L}\right)+b_n\sin\left(\frac{n\pi t}{L}\right).
\]
Die Koeffizienten beschreiben, wie stark die Frequenz $n\pi/L$ in $f$ vorkommt.

\subsection{1.8 Signalübertragung}
Statt das ganze Signal zu senden, kann man die ersten paar Koeffizienten übertragen. Das funktioniert, wenn die Reihe schnell konvergiert.

\subsection{1.9 Approximation}
Da die Reihe schnell konvergiert, kann man damit Funktionen gut approximieren. Man wählt den Definitionsbereich als Periode und kann so auch nichtperiodische Funktionen periodisch approximieren.

\subsection{2.1 Differentialgleichungen lösen}
Man kann Differentialgleichungen lösen, indem man etwa für $\cos t$ die allgemeine Lösung als Fourierreihe ansetzt:
\[
\sum a_n\cos(nt)
\]
und Koeffizienten vergleicht.

% =========================================================
\section{Woche 12: Fourier-Analysis}

\subsection{Einführung}
In der linearen Algebra schreibt man einen Vektor als Linearkombination von Basisvektoren:
\[
\binom34=a\binom10+b\binom01.
\]
Die Koeffizienten findet man über das Skalarprodukt.

Statt Vektoren betrachten wir jetzt Funktionen. Funktionen bestehen aus Linearkombinationen von Basisfunktionen.

\subsection{1.1 Fourierreihe}
Sei $f(t)$ $2L$-periodisch. Dann
\begin{importantbox}
\[
f(t)=\frac{a_0}{2}+\sum_{n=1}^{\infty}a_n\cos\left(\frac{n\pi}{L}t\right)+b_n\sin\left(\frac{n\pi}{L}t\right).
\]
\end{importantbox}
Basisfunktionen sind $\cos(n\pi t/L)$ und $\sin(n\pi t/L)$. Die Koeffizienten $a_n,b_n$ beschreiben, wie viel von der Frequenz $n\pi/L$ in $f$ steckt.

\subsection{1.2 Skalarprodukt und Koeffizienten}
Für Funktionen lautet das Skalarprodukt
\[
\langle f,g\rangle=\int_{-L}^{L}f(t)g(t)\dd t.
\]
Wie bei Vektoren nimmt man das Skalarprodukt mit den Basisfunktionen:
\begin{importantbox}
\[
a_n=\frac1L\int_{-L}^{L}f(t)\cos\left(\frac{n\pi}{L}t\right)\dd t,
\]
\[
b_n=\frac1L\int_{-L}^{L}f(t)\sin\left(\frac{n\pi}{L}t\right)\dd t.
\]
\end{importantbox}

\subsection{1.3 Komplexe Fourierreihe}
Wenn man negative $n$ zulässt, kann man die Fourierreihe mit einem Koeffizienten ausdrücken:
\begin{importantbox}
\[
f(t)=\sum_{n=-\infty}^{\infty}c_n\ee^{\ii n\pi t/L},
\qquad
c_n=\frac1{2L}\int_{-L}^{L}f(t)\ee^{-\ii n\pi t/L}\dd t.
\]
\end{importantbox}
Die Darstellungen sind verbunden durch
\[
c_n=\frac12(a_n-\ii b_n),
\qquad
c_{-n}=\frac12(a_n+\ii b_n).
\]

\subsection{1.4 Gerade und ungerade Funktionen}
\begin{fragebox}
Falls $f$ gerade oder ungerade ist, was lässt sich über $a_n$ und $b_n$ sagen?
\end{fragebox}
\begin{antwortbox}
Da $\cos$ gerade und $\sin$ ungerade ist:
\begin{itemize}
\item Ist $f$ ungerade, dann $a_n=0$.
\item Ist $f$ gerade, dann $b_n=0$.
\end{itemize}
\end{antwortbox}

\subsection{1.5 Periodische Fortsetzungen}
Fourierreihen gelten für periodische Funktionen. Falls $f$ nicht periodisch, aber beschränkt auf einem Intervall ist, setzt man $f$ periodisch fort.

\begin{center}
\begin{tikzpicture}[scale=0.85]
  \draw[axis] (-0.5,0)--(4,0) node[right] {$t$};
  \draw[axis] (0,-0.3)--(0,1.8) node[above] {$f(t)$};
  \draw[blue,thick] (0,1)--(1,1);
  \node at (0.5,-0.25) {$[0,L]$};
\end{tikzpicture}
\end{center}
Es gibt drei Arten:
\begin{enumerate}[label=\Roman*.]
\item $L$-periodische Fortsetzung.
\item Gerade $2L$-periodische Fortsetzung.
\item Ungerade $2L$-periodische Fortsetzung.
\end{enumerate}

\begin{center}
\begin{tikzpicture}[scale=0.75]
  \begin{scope}
    \draw[axis] (-2.5,0)--(3,0) node[right] {$t$}; \draw[axis] (0,-0.3)--(0,1.6) node[above] {$f$};
    \foreach \k in {-2,-1,0,1,2}{\draw[blue,thick] (\k,1)--(\k+0.7,1);}
    \node at (0.2,-0.55) {I: $L$-periodisch};
  \end{scope}
  \begin{scope}[xshift=7cm]
    \draw[axis] (-2.5,0)--(3,0) node[right] {$t$}; \draw[axis] (0,-0.3)--(0,1.6) node[above] {$f$};
    \draw[blue,thick] (-2,1)--(-1.2,1); \draw[blue,thick] (-0.8,1)--(0.8,1); \draw[blue,thick] (1.2,1)--(2,1);
    \node at (0.2,-0.55) {II: gerade};
  \end{scope}
\end{tikzpicture}
\end{center}

\subsection{1.6 Fortsetzungen zeichnen}
\begin{fragebox}
Zeichne die drei Fortsetzungen einer auf $[0,L]$ linearen Funktion.
\end{fragebox}
\begin{antwortbox}
Die drei Skizzen entsprechen: periodische Sägezahnfortsetzung, gerade Dreiecksfortsetzung und ungerade Sägezahnfortsetzung.
\end{antwortbox}

\begin{center}
\begin{tikzpicture}[scale=0.7]
  \begin{scope}
    \draw[axis] (-2.2,0)--(3,0); \draw[axis] (0,-0.2)--(0,1.7) node[above] {$f$};
    \foreach \k in {-2,-1,0,1,2}{\draw[blue,thick] (\k,0)--(\k+1,1);}
    \node at (0.5,-0.5) {periodisch};
  \end{scope}
  \begin{scope}[xshift=6cm]
    \draw[axis] (-2.2,0)--(3,0); \draw[axis] (0,-0.2)--(0,1.7) node[above] {$f$};
    \draw[blue,thick] (-2,0)--(-1,1)--(0,0)--(1,1)--(2,0);
    \node at (0.5,-0.5) {gerade};
  \end{scope}
  \begin{scope}[yshift=-3.1cm]
    \draw[axis] (-2.2,0)--(3,0); \draw[axis] (0,-1.3)--(0,1.3) node[above] {$f$};
    \foreach \k in {-2,-1,0,1,2}{\draw[blue,thick] (\k,0)--(\k+1,1); \draw[blue,thick] (\k,-1)--(\k+1,0);}
    \node at (0.5,-1.6) {ungerade};
  \end{scope}
\end{tikzpicture}
\end{center}

\subsection{1.7 Beispiel Koeffizienten berechnen}
\begin{beispielbox}
Sei $f(x)=x^2+x+2$ auf $[-3\pi,3\pi]$ definiert und mit Periode $6\pi$ fortgesetzt. Bestimme $w$, $a_0$, $a_n$, $b_n$.
\end{beispielbox}
\begin{tippbox}
Nutze Symmetrien der Integrale: Produkte mit $\sin$ bzw. $\cos$ können auf geraden und ungeraden Anteilen vereinfacht werden.
\end{tippbox}
Hier ist
\[
w=\frac{2\pi}{T}=\frac{2\pi}{6\pi}=\frac13.
\]
Außerdem
\[
a_0=\frac1L\int_{-L}^{L}f(x)\dd x,
\qquad L=3\pi.
\]
Aus der Rechnung ergibt sich
\[
a_0=6\pi^2+4.
\]
Weiter
\[
a_n=\frac1{3\pi}\int_{-3\pi}^{3\pi}(x^2+x+2)\cos\left(\frac{nx}{3}\right)
\dd x
=\frac{36(-1)^n}{n^2},
\]
und
\[
b_n=\frac1{3\pi}\int_{-3\pi}^{3\pi}(x^2+x+2)\sin\left(\frac{nx}{3}\right)
\dd x
=\frac{6(-1)^{n+1}}{n}.
\]

\subsection{1.8 Diskrete Fouriertransformation}
Sei
\[
f(t)=\sum_{n=-\infty}^{\infty}c_n\ee^{\ii n\pi t/L}.
\]
Dann sind die $c_n$ die diskrete Fouriertransformation. Man liegt nicht mehr im Zeitbereich, sondern im Frequenzbereich.

\begin{importantbox}
$f(t)$ fragt: Wie groß ist $f$ zu einer gegebenen Zeit? \\
$\widehat f(\omega)$ fragt: Wie groß ist $f$ bei einer gegebenen Frequenz?
\end{importantbox}

\subsection{1.9 Fouriertransformation}
Falls $f$ nicht beschränkt ist oder man keine periodische Fortsetzung möchte, nimmt man die Periode im Grenzfall gegen $\infty$. Aus der Summe wird ein Integral und aus diskreten Frequenzen ein kontinuierliches Frequenzspektrum.

\begin{importantbox}
\[
f(t)=\frac1{\sqrt{2\pi}}\int_{-\infty}^{\infty}\widehat f(\omega)\ee^{\ii\omega t}\dd\omega,
\qquad
\widehat f(\omega)=\frac1{\sqrt{2\pi}}\int_{-\infty}^{\infty}f(t)\ee^{-\ii\omega t}\dd t.
\]
\end{importantbox}

% =========================================================
\section{Woche 13: Fouriertransformation}

\subsection{1.1 Existenz}
\begin{fragebox}
Wann existiert eine Fouriertransformation?
\end{fragebox}
\begin{antwortbox}
Das Integral darf nicht gegen unendlich gehen. Die Funktion muss absolut integrierbar sein:
\[
\int_{-\infty}^{\infty}|f(t)|\dd t<\infty.
\]
\end{antwortbox}

\subsection{1.2 Beispiele zur absoluten Integrierbarkeit}
\begin{beispielbox}
\begin{itemize}
\item $f(t)=\ee^t$ ist nicht absolut integrierbar.
\item $f(t)=\mathbf 1_{[-1,1]}(t)$ ist absolut integrierbar.
\item $f(t)=\ee^{-t}\mathbf 1_{t\ge0}$ ist absolut integrierbar.
\end{itemize}
\end{beispielbox}

\subsection{1.3 Fouriertransformation von $\ee^{-t}\mathbf 1_{t\ge0}$}
\begin{beispielbox}
Berechne die Fouriertransformation von
\[
f(t)=\begin{cases}\ee^{-t},&t\ge0,\\0,&\text{sonst.}\end{cases}
\]
\end{beispielbox}
\[
\widehat f(\omega)=\frac1{\sqrt{2\pi}}\int_0^{\infty}\ee^{-t}\ee^{-\ii\omega t}\dd t
=\frac1{\sqrt{2\pi}}\int_0^{\infty}\ee^{-(1+\ii\omega)t}\dd t.
\]
Da $\Ree(1+\ii\omega)=1>0$, ist
\[
\widehat f(\omega)=\frac1{\sqrt{2\pi}}\frac1{1+\ii\omega}.
\]

\subsection{1.4 Allgemeines Beispiel}
Für
\[
f(t)=\begin{cases}\ee^{-at},&t\ge0,\\0,&\text{sonst,}\end{cases}
\qquad a>0,
\]
gilt
\begin{importantbox}
\[
\widehat f(\omega)=\frac1{\sqrt{2\pi}}\frac1{a+\ii\omega}.
\]
\end{importantbox}

\subsection{1.5 Eigenschaften der Fouriertransformation}
Sei $\mathcal F[f]=\widehat f$. Wichtige Rechenregeln sind:
\begin{importantbox}
\begin{enumerate}[label=\arabic*.]
\item Verschiebung:
\[
\mathcal F[f(t-a)](\omega)=\ee^{-\ii\omega a}\widehat f(\omega).
\]
\item Modulation:
\[
\mathcal F[\ee^{\ii a t}f(t)](\omega)=\widehat f(\omega-a).
\]
\item Skalierung:
\[
\mathcal F[f(at)](\omega)=\frac1{|a|}\widehat f\left(\frac{\omega}{a}\right).
\]
\item Doppelte Fouriertransformation:
\[
\mathcal F(\mathcal F(f))=f(-t).
\]
\item Ableitungen:
\[
\mathcal F[f'](\omega)=\ii\omega\widehat f(\omega),
\qquad
\mathcal F[t^nf(t)](\omega)=\ii^n\frac{d^n}{d\omega^n}\widehat f(\omega).
\]
\item Wichtig für Verschiebung und Skalierung:
\[
\mathcal F[f(a(t-b))](\omega)=\frac1{|a|}\ee^{-\ii\omega b}\widehat f\left(\frac{\omega}{a}\right).
\]
\end{enumerate}
\end{importantbox}

\subsection{1.6 Beispiel mit Verschiebung}
\begin{beispielbox}
Sei $f(t)=\ee^{-|t|}$ und
\[
\widehat f(\omega)=\frac{2}{1+\omega^2}.
\]
Was ist $\mathcal F[\ee^{-|t-2|}]$?
\end{beispielbox}
Da $\ee^{-|t-2|}=f(t-2)$, verwendet man die Verschiebungsregel:
\[
\mathcal F[f(t-2)](\omega)=\ee^{-2\ii\omega}\frac{2}{1+\omega^2}.
\]

\subsection{1.7 Frage}
\begin{fragebox}
Sei $f(t)=\ee^{-|t|}$ und $\widehat f(\omega)=\frac{2}{1+\omega^2}$. Was ist
\[
\mathcal F\left[\ee^{-\left|\frac{t+3}{2}\right|}\right]?
\]
\end{fragebox}
\begin{antwortbox}
Man schreibt
\[
\ee^{-\left|\frac{t+3}{2}\right|}=f\left(\frac{t+3}{2}\right)=f\left(\frac12(t+3)\right).
\]
Mit Skalierung und Verschiebung erhält man
\[
\mathcal F\left[f\left(\frac{t+3}{2}\right)\right](\omega)
=2\ee^{3\ii\omega}\widehat f(2\omega)
=2\ee^{3\ii\omega}\frac{2}{1+4\omega^2}.
\]
\end{antwortbox}

\subsection{1.8 Faltung}
\begin{importantbox}
Die Faltung ist definiert durch
\[
(f*g)(t)=\int_{-\infty}^{\infty}f(t-x)g(x)\dd x
=\int_{-\infty}^{\infty}f(x)g(t-x)\dd x.
\]
\end{importantbox}

\begin{beispielbox}
Sei
\[
f(t)=\begin{cases}1,&t\ge0,\\0,&\text{sonst.}\end{cases}
\]
Berechne $(\ee^{-3t}f(t))*f(t)$.
\end{beispielbox}
Für $t\ge0$:
\[
(\ee^{-3t}f(t))*f(t)=\int_0^t \ee^{-3x}\dd x=\left[-\frac13\ee^{-3x}\right]_0^t=\frac{1-\ee^{-3t}}{3}.
\]
Für $t<0$ ist der Wert $0$. Also
\[
(\ee^{-3t}f(t))*f(t)=\frac{1-\ee^{-3t}}{3}f(t).
\]
Allgemein:
\begin{importantbox}
\[
(\ee^{-at}f(t))*f(t)=\frac{1-\ee^{-at}}{a}f(t).
\]
\end{importantbox}

% =========================================================
\section{Woche 14: Faltung und Laplace-Transformation}

\subsection{1.1 Tiefpassfilter}
Man möchte ein Zeitsignal $f(t)$ durch einen Tiefpassfilter laufen lassen. Im Frequenzbereich multipliziert man die Fouriertransformation $\mathcal F[f]$ mit
\[
\chi_{[-1,1]}(\omega).
\]
Hohe Frequenzen werden abgeschnitten:
\[
\widehat T(\omega)=\widehat f(\omega)\chi_{[-1,1]}(\omega).
\]
Im Zeitbereich entspricht das einer Faltung:
\[
T(t)=f*\mathcal F^{-1}[\chi_{[-1,1]}](t).
\]
Da
\[
\mathcal F^{-1}[\chi_{[-1,1]}](t)=\sqrt{\frac2\pi}\frac{\sin t}{t},
\]
erhält man
\[
T(t)=f(t)*\sqrt{\frac2\pi}\frac{\sin t}{t}.
\]

\begin{center}
\begin{tikzpicture}[scale=0.8]
  \draw[axis] (-3,0)--(3,0) node[right] {$\omega$};
  \draw[axis] (0,-0.2)--(0,1.6) node[above] {$\chi_{[-1,1]}$};
  \draw[blue,thick] (-1,0)--(-1,1)--(1,1)--(1,0);
  \node[below] at (-1,0) {$-1$}; \node[below] at (1,0) {$1$};
\end{tikzpicture}
\end{center}

\subsection{1.2 Faltungen in KI}
Faltungen werden in der KI ständig verwendet. Convolutional Neural Networks benutzen Faltungen als Filter, um Features aus einem Datensatz zu extrahieren.

\subsection{1.3 Differentialgleichung mit Fouriertransformation}
\begin{fragebox}
Berechne $\widehat y(\omega)$, wobei
\[
-y''(t)+y(t)=\chi_{[-1,1]}(t).
\]
\end{fragebox}
\begin{antwortbox}
Fouriertransformieren liefert
\[
\omega^2\widehat y(\omega)+\widehat y(\omega)=\mathcal F[\chi_{[-1,1]}](\omega).
\]
Da
\[
\mathcal F[\chi_{[-1,1]}](\omega)=\sqrt{\frac2\pi}\frac{\sin\omega}{\omega},
\]
folgt
\[
\widehat y(\omega)=\frac{\sqrt{2/\pi}\,\sin\omega}{\omega(1+
\omega^2)}.
\]
\end{antwortbox}

\subsection{1.4 Von Fourier zu Laplace}
Die Fouriertransformation ist
\[
\widehat f(\omega)=\frac1{\sqrt{2\pi}}\int_{-\infty}^{\infty}f(t)\ee^{-\ii\omega t}\dd t.
\]
Hier ist $\omega$ reell. Ersetzt man $\omega$ durch eine komplexe Zahl $s=a+\ii b$, erhält man Dämpfungsfaktoren. Für $a>0$ wächst $\ee^{-st}$ bei negativen $t$, für $a<0$ bei positiven $t$. Deshalb integriert man nur über eine Seite, üblicherweise von $0$ bis $\infty$.

\begin{importantbox}
Die Laplace-Transformation ist
\[
\mathcal L[f](s)=\int_0^{\infty}f(t)\ee^{-st}\dd t.
\]
\end{importantbox}

\subsection{1.5 Idee der Laplace-Transformation}
Die Laplace-Transformation ist im Kern eine Fouriertransformation von $f(t)\ee^{-at}$. Man dämpft also $f$, bevor man transformiert. Dadurch kann man viel mehr Funktionen transformieren: $f$ darf höchstens exponentiell wachsen.

Wenn
\[
|f(t)|\le Me^{a_0t},
\]
dann muss
\[
\Ree(s)>a_0
\]
gelten, damit genügend gedämpft wird.

\subsection{1.6 Beispiel: Laplace von $\ee^{bt}H(t)$}
\begin{beispielbox}
Berechne $\mathcal L[f](s)$ für $f(t)=\ee^{bt}H(t)$.
\end{beispielbox}
\[
\mathcal L[f](s)=\int_0^{\infty}\ee^{bt}\ee^{-st}\dd t
=\int_0^{\infty}\ee^{-(s-b)t}\dd t
=\frac1{s-b},
\qquad \Ree(s)>\Ree(b).
\]

\subsection{1.7 Nicht transformierbare Funktion}
$\ee^{t^2}$ kann mit Laplace nicht transformiert werden, weil es schneller als exponentiell wächst.

\subsection{1.8 Beispiel: Rücktransformation}
\begin{beispielbox}
Sei
\[
f(t)=3+10t-3\cos t+2\sin t.
\]
Verwende die Laplace-Tabelle.
\end{beispielbox}
Mit
\[
\mathcal L[1]=\frac1s,
\quad
\mathcal L[t]=\frac1{s^2},
\quad
\mathcal L[\cos(at)]=\frac{s}{s^2+a^2},
\quad
\mathcal L[\sin(at)]=\frac{a}{s^2+a^2},
\]
ergibt sich
\[
\mathcal L[f](s)=\frac3s+\frac{10}{s^2}-\frac{3s}{s^2+1}+\frac{2}{s^2+1}.
\]
Für rationale Funktionen kann man Partialbruchzerlegung verwenden, um die Rücktransformation abzulesen.

\subsection{1.9 Wichtige Formeln für Laplace}
Die folgende Tabelle fasst die wichtigsten Transformationen zusammen.

\begin{longtable}{>{\(}l<{\)}>{\(}l<{\)}}
\toprule
f(t) & \mathcal L[f](s) \\
\midrule
\endhead
1 & 1/s \\
t^n & n!/s^{n+1} \\
\sin(at) & a/(s^2+a^2) \\
\cos(at) & s/(s^2+a^2) \\
\ee^{at} & 1/(s-a) \\
\ee^{at}\sin(bt) & b/((s-a)^2+b^2) \\
\ee^{at}\cos(bt) & (s-a)/((s-a)^2+b^2) \\
t^n\ee^{at} & n!/(s-a)^{n+1} \\
af(t)+bg(t) & a\mathcal L[f](s)+b\mathcal L[g](s) \\
tf(t) & -\frac{d}{ds}\mathcal L[f](s) \\
t^nf(t) & (-1)^n\frac{d^n}{ds^n}\mathcal L[f](s) \\
f'(t) & s\mathcal L[f](s)-f(0) \\
f''(t) & s^2\mathcal L[f](s)-sf(0)-f'(0) \\
\ee^{at}f(t) & \mathcal L[f](s-a) \\
\int_0^t f(\tau)g(t-\tau)\dd\tau & \mathcal L[f](s)\mathcal L[g](s) \\
\bottomrule
\end{longtable}

\end{document}
